\documentclass[11pt]{article}
\usepackage[utf8]{inputenc}
\usepackage[english]{babel}
\usepackage{geometry}
\usepackage{amsmath}
\usepackage{booktabs}
\usepackage{xcolor}
\usepackage{mdframed}

\geometry{a4paper, margin=1in}

\title{Technical Specifications: USRP B200mini}
\author{Radiofrequency Engineering}
\date{\today}

% Warning box
\newmdenv[
  backgroundcolor=red!8,
  linecolor=red!60,
  linewidth=1pt,
  roundcorner=4pt,
  innertopmargin=6pt,
  innerbottommargin=6pt
]{warningbox}

\begin{document}

\maketitle
\tableofcontents
\newpage

% =========================================================
\section{Gains}
% =========================================================
The USRP B200mini uses the AD9364 transceiver. Through the UHD driver, the multiple physical gain stages are abstracted into a single programmable gain-amplification parameter (PGA).

\begin{itemize}
    \item \textbf{RX Gain (PGA):} Range from \textbf{0.0 dB to 76.0 dB}. Controls the receive chain, including the LNA and baseband stages.
    \item \textbf{TX Gain (PGA):} Range from \textbf{0.0 dB to 89.8 dB}. Adjusts the RF signal output power.
    \item \textbf{AGC Modes (RX):}
    \begin{itemize}
        \item \textit{Manual:} The user defines a fixed gain value.
        \item \textit{Fast Attack:} Fast automatic adjustment for burst-type signals.
        \item \textit{Slow Attack:} Slow automatic adjustment for constant-power signals.
    \end{itemize}
\end{itemize}

\textbf{Operating note:} It is recommended to use at least half of the available gain to obtain a reasonable dynamic range.

% =========================================================
\section{Hardware Filters}
% =========================================================
The USRP B200mini combines analog baseband filters with digital filters automatically configured by UHD.

\begin{enumerate}
    \item \textbf{Analog filters (Baseband LPF)}

    \textbf{Modifiable through:} \texttt{set\_rx\_bandwidth(bw)} and \texttt{set\_tx\_bandwidth(bw)}.

    \textbf{What it does:} Directly adjusts the Butterworth low-pass filter. Any value between 200 kHz and 56 MHz can be set. If it is not configured manually, UHD attempts to adjust it automatically according to the \texttt{sample\_rate} value.

    \item \textbf{Digital filters (Half-Band and FIR)}

    \textbf{Modifiable through:} They are not modified directly. They are configured implicitly when using \texttt{set\_sample\_rate(rate)} and \texttt{set\_master\_clock\_rate(rate)}.

    \textbf{What it does:} UHD automatically calculates the required interpolation or decimation stages and configures the Half-Band filters and digital FIR coefficients to prevent aliasing. Everything happens in the background.
\end{enumerate}

In code, it is enough to set the bandwidth (\texttt{bandwidth}) and sample rate (\texttt{sample\_rate}); UHD performs the required calculations and programs the AD9364 transceiver.

% =========================================================
\section{Frequency Ranges}
% =========================================================
Nominal operating range: \textbf{70 MHz to 6 GHz}.

% =========================================================
\section{Sample Rate}
% =========================================================
Maximum sample rate: \textbf{61.44 MS/s} (single-channel SDR).

% =========================================================
\section{Bits (Resolution)}
% =========================================================
\textbf{12-bit} ADC and DAC resolution.

% =========================================================
\section{IQ \& DC Corrections}
% =========================================================
The USRP B200mini already has native I/Q imbalance and DC leakage (\textit{DC offset} or \textit{DC spike}) corrections implemented. The system handles them at different levels:

\begin{itemize}
    \item \textbf{At the hardware level (AD9364 transceiver):} The chip integrates tracking algorithms that continuously correct \textit{DC offset}, both in the RF stages and in baseband, and they are enabled by default. Additionally, the transceiver has integrated I/Q correction through matrix multipliers.
    \item \textbf{At the software level (UHD driver):} The UHD abstraction enables and manages these automatic factory corrections, so it is not necessary to program corrective filters from scratch in user code.
    \item \textbf{Specific calibration:} If automatic adjustment is not sufficient, UHD includes native command-line utilities, such as \texttt{uhd\_cal\_tx\_dc\_offset} or \texttt{uhd\_cal\_rx\_iq\_balance}, which allow test signals to be injected and static calibration tables stored in the hardware to fine-tune the values.
\end{itemize}

In summary, these corrections work \textit{out of the box}, but manual calibrations can be forced through the driver if the application requires it.

% =========================================================
\section{SDR Inputs and Connectors}
% =========================================================
The USRP B200mini is a 1x1 full-duplex SDR. Its main connections are concentrated in the USB 3.0 interface, the RF ports, and the external reference input for synchronization.

\begin{table}[h]
    \centering
    \begin{tabular}{p{0.16\textwidth}p{0.26\textwidth}p{0.46\textwidth}}
        \toprule
        \textbf{Connector} & \textbf{Function} & \textbf{Main notes} \\
        \midrule
        USB3 & Power and data & USB 3.0 interface to the host computer. \\
        J1 TRX & TX/RX RF & Shared TX/RX port; TX output up to +20 dBm. \\
        J2 RX2 & RF reception & Dedicated receive input; maximum input power -15 dBm. \\
        J3 Ref/PPS & External reference & Input for 10 MHz reference or PPS pulse; +15 dBm max. \\
        J5 JTAG & FPGA JTAG & Access to the Spartan-6 FPGA for development and diagnostics. \\
        J6 GPIO & Digital control & GPIO pins connected to the FPGA for auxiliary signals. \\
        \bottomrule
    \end{tabular}
\end{table}

\begin{itemize}
    \item \textbf{Internal clock:} The master clock feeds the RF and DSP chains. UHD can select it automatically according to the requested sample rate, or the user can fix it through \texttt{master\_clock\_rate}. The typical configurable range is 5 MHz to 61.44 MHz.
    \item \textbf{External Reference Clock:} The J3 input allows the SDR to be disciplined with an external 10 MHz reference. This is useful when frequency coherence between equipment or synchronization with a laboratory source is required.
    \item \textbf{External PPS:} The same J3 connector can be used as a PPS input to align the equipment time reference. In multi-SDR configurations, it is recommended to explicitly fix the \texttt{master\_clock\_rate} to avoid automatic changes that break synchronization.
    \item \textbf{RX port:} J2 RX2 is the dedicated receive input. Applying more than -15 dBm should be avoided to protect the RF front end.
    \item \textbf{TX/RX port:} J1 TRX works as a shared transmit and receive port. During transmission it must be connected to an antenna or 50 $\Omega$ load; in loopback tests, using at least 30 dB of attenuation is recommended.
\end{itemize}

% =========================================================
\section{Safe Operating Limits}
% =========================================================

\begin{warningbox}
\textbf{Critical warnings - read before connecting the equipment}
\end{warningbox}

\vspace{4pt}
\begin{itemize}
    \item \textbf{Maximum RX power:} Never apply more than \textbf{-15 dBm} to any RF input (J1 TRX or J2 RX2). Exceeding this level can permanently damage the front end.
    \item \textbf{Loopback:} In TX$\to$RX loopback configuration, use \textbf{at least 30 dB of attenuation} between the ports to protect the input.
    \item \textbf{TX port without antenna:} The J1 TRX port must always be terminated with an antenna or a \textbf{50 $\Omega$} load when transmitting. Transmitting with the port open can damage the amplifier.
    \item \textbf{External reference J3:} 10 MHz reference signal with a maximum level of \textbf{+15 dBm}.
    \item \textbf{ESD and handling:} Always handle with antistatic protection. Do not allow metal objects to touch the PCB while it is powered.
    \item \textbf{Maximum TX power:} The output of the J1 TRX port can reach \textbf{+20 dBm} (peak, frequency-dependent); the typical guaranteed value is $>$10 dBm.
\end{itemize}

% =========================================================
\section{Master Clock Rate - Rules and Pitfalls}
% =========================================================

The \texttt{master\_clock\_rate} (MCR) is the most critical configuration parameter for extracting maximum performance from the device and the main source of silent problems.

\subsection*{Range and valid values}
\begin{itemize}
    \item Configurable range: \textbf{5 MHz to 61.44 MHz}.
    \item Values above 56 MHz are possible but \textbf{not recommended}.
    \item The MCR must be an \textbf{integer multiple of the sample rate}. If this condition is not met, UHD displays a warning and internally applies a rate different from the requested one.
\end{itemize}

\subsection*{Automatic vs.\ manual behavior}
By default, UHD selects the MCR automatically according to the requested \texttt{sample\_rate}, choosing the highest possible value to enable the largest number of Half-Band filters. Automatic mode is disabled as soon as \texttt{set\_master\_clock\_rate()} is called or \texttt{master\_clock\_rate} is passed as an initialization argument.

\begin{verbatim}
# Set MCR manually (Python UHD)
usrp.set_master_clock_rate(32e6)

# Or through a device argument
uhd_usrp_probe --args="master_clock_rate=52e6"
\end{verbatim}

\subsection*{Known pitfalls}
\begin{itemize}
    \item \textbf{MCR change breaks synchronization:} Any MCR change reconfigures the entire clock chain, invalidating previous time configurations (PPS, multi-USRP sync).
    \item \textbf{Multi-USRP:} When synchronizing multiple units, \textbf{fix the MCR explicitly} on all of them. Do not depend on automatic mode, since it can select different values per unit.
    \item \textbf{Locked MCR prevents sample-rate changes:} If the MCR is fixed and a \texttt{sample\_rate} that is not an integer submultiple is requested, UHD will apply the nearest valid rate without making it clearly visible.
\end{itemize}

% =========================================================
\section{Indicator LEDs}
% =========================================================

The LEDs allow device status to be diagnosed in the field without needing a console.

\begin{table}[h]
    \centering
    \begin{tabular}{p{0.14\textwidth}p{0.22\textwidth}p{0.52\textwidth}}
        \toprule
        \textbf{LED} & \textbf{Function} & \textbf{States} \\
        \midrule
        PWR & Power & Off = no power \quad On = powered (USB or external) \\
        \addlinespace
        TRX & TX/RX activity & Off = no activity \quad Green = receiving \quad Red = transmitting \quad Orange = alternating TX/RX \\
        \addlinespace
        RX2 & RX2 activity & Off = no activity \quad Green = receiving \\
        \addlinespace
        S0 & Reference lock & Off = no lock \quad Green = reference locked \\
        \addlinespace
        S1 & Reference presence & Off = low or absent signal \quad Green = reference signal present \\
        \bottomrule
    \end{tabular}
\end{table}

\textbf{Interpretation of blinks on TRX/RX2:}
\begin{itemize}
    \item A blink on the TX LED indicates \textbf{underflow}: the host is not delivering samples to the SDR on time. It is reported in the console as ``\texttt{U}''.
    \item A blink on the RX LED indicates \textbf{overflow}: the host is not consuming samples from the SDR on time. It is reported as ``\texttt{O}''.
    \item If this occurs, reduce the \texttt{sample\_rate}, increase the system buffer, or review CPU/USB load.
\end{itemize}

% =========================================================
\section{GPIO Pinout (J6)}
% =========================================================

Connector J6 exposes 8 GPIO lines connected directly to the Spartan-6 FPGA, operating at \textbf{3.3 V}. It is useful for control signals, external trigger, or integration with auxiliary hardware.

\begin{table}[h]
    \centering
    \begin{tabular}{ccc|ccc}
        \toprule
        \textbf{Pin} & \textbf{Signal} & \textbf{Level} & \textbf{Pin} & \textbf{Signal} & \textbf{Level} \\
        \midrule
        1 & 3.3 V (VCC) & Power & 7  & 3.3 V (VCC) & Power \\
        2 & GPIO\_0      & 3.3 V logic & 8  & GPIO\_4      & 3.3 V logic \\
        3 & GPIO\_1      & 3.3 V logic & 9  & GPIO\_5      & 3.3 V logic \\
        4 & GPIO\_2      & 3.3 V logic & 10 & GPIO\_6      & 3.3 V logic \\
        5 & GPIO\_3      & 3.3 V logic & 11 & GPIO\_7      & 3.3 V logic \\
        6 & GND          & Reference   & 12 & GND          & Reference \\
        \bottomrule
    \end{tabular}
\end{table}

\textbf{Usage notes:}
\begin{itemize}
    \item GPIO\_0 to GPIO\_7 lines are configurable as input or output from UHD through the GPIO API of \texttt{multi\_usrp}.
    \item Direction and value control is performed with \texttt{set\_gpio\_attr()} by specifying bank \texttt{"FP0"}.
    \item For applications that require UART signals over GPIO, a custom FPGA bitfile must be compiled.
\end{itemize}

\end{document}
